% =====================================================================
% 1. Introdução
% =====================================================================
\section{Introdução}

Boa parte dos modelos usados em engenharia nasce contínua: tensões, correntes, campos e envelopes de sinais são descritos por equações diferenciais. O computador, por sua vez, só executa um número finito de operações aritméticas sobre quantidades representadas com precisão finita. O cálculo numérico é a disciplina que torna esse encontro rigoroso. Ele não substitui a modelagem; ele diz em que condições uma sequência discreta de operações produz uma aproximação confiável do modelo, quanto isso custa e quando o resultado deixa de ter sentido físico.

A história dessa disciplina é mais antiga do que a do computador eletrônico. Tabelas, logaritmos, interpolação e os primeiros métodos para equações diferenciais foram criados para o cálculo manual e para a produção de tábuas matemáticas. As máquinas de Pascal, de Leibniz e de Babbage, e depois os computadores programáveis do século XX, automatizaram procedimentos que a análise numérica já havia formulado. Compreender essa trajetória evita tratar o algoritmo como uma caixa atual, desligada das limitações de arredondamento, de estabilidade e de custo que o produziram.

O tema específico deste trabalho é a integração numérica de equações diferenciais ordinárias por métodos de passo único, com aprofundamento no método de Euler, no Euler melhorado e no Runge-Kutta clássico de quarta ordem. O problema de engenharia que conduz os experimentos é o circuito RLC série. Com uma escolha de parâmetros, o mesmo circuito descreve o transitório de um estágio ressonante: o tempo de acomodação de um filtro, o crescimento da corrente quando uma portadora entra em sintonia e o toque que um pulso curto deixa num tanque seletivo. Há solução analítica em vários casos, o que permite medir erro de verdade, e não apenas a diferença entre dois métodos.

O objetivo geral é elaborar um estudo que una a evolução histórica do cálculo numérico ao desenvolvimento da computação e que implemente, teste e critique um método numérico de interesse direto para a engenharia de computação e de telecomunicações. Os objetivos específicos são: situar o cálculo aproximado desde as tabelas até as bibliotecas científicas atuais; relacionar essa evolução aos instrumentos de cálculo e aos computadores eletrônicos; discutir ponto flutuante, erro, estabilidade e custo na implementação; aplicar os métodos a um circuito ressonante; e registrar, com código próprio, a convergência, a estabilidade e o efeito da precisão finita.

O texto está organizado da seguinte forma. A Seção~\ref{sec:historia} percorre a história pedida, das tábuas matemáticas à computação científica contemporânea. A Seção~\ref{sec:teoria} formula o problema de valor inicial, o modelo do circuito e a análise de erro e de estabilidade dos três métodos. A Seção~\ref{sec:algoritmos} descreve os algoritmos e a implementação. A Seção~\ref{sec:aplicacoes} liga o circuito às grandezas usadas em comunicações. A Seção~\ref{sec:resultados} discute os experimentos. A Seção~\ref{sec:conclusao} sintetiza os resultados e indica continuações. O código completo está no Apêndice~\ref{ap:codigo}.
